\documentclass[11pt,a4paper]{article}
\usepackage[utf8]{inputenc}
\usepackage[margin=1in]{geometry}
\usepackage{amsmath}
\usepackage{amssymb}
\usepackage{enumitem}
\usepackage{booktabs}
\usepackage{titlesec}
\usepackage{microtype}

\titleformat{\section}{\large\bfseries\sffamily}{}{0em}{}[\titlerule]
\titleformat{\subsection}{\normalsize\bfseries\sffamily}{}{0em}{}

\setlength{\parindent}{0pt}
\setlength{\parskip}{6pt plus 2pt minus 1pt}

\title{\textbf{Hard CAMS Exam Questions: Missing Topics Edition}}
\author{\textbf{Advanced AML Training}}
\date{\today}

\begin{document}
	
	\maketitle
	
	\section{Terrorist Financing (TF) Typologies}
	
	\subsection{1. Funding Sources for Terrorism}
	A Financial Intelligence Unit analysis reveals that a domestic terrorism cell received multiple small-dollar wire transfers ($200–$400 each) from over 50 distinct individuals across three different countries. The senders have no apparent connection to each other or to the cell. No single transaction exceeds local reporting thresholds. What is the most accurate classification of this financing method?
	\begin{enumerate}[label=\Alph*.]
		\item Trade-based money laundering
		\item Crowdfunding-enabled terrorist financing (micro-donations)
		\item Black Market Peso Exchange adaptation
		\item Structuring under FATF Recommendation 14
	\end{enumerate}
	\textbf{Correct Answer: B}
	
	\subsection{2. NPO Abuse for TF}
	A non-profit organization registered in a low-risk jurisdiction receives a large donation from a foreign charity operating in a conflict zone. The local NPO then distributes funds to what appear to be community development projects, but transaction monitoring shows that 80\% of the funds are withdrawn in cash within 48 hours of receipt at a branch located 300km from the NPO's registered address. What red flag is most indicative of terrorist financing?
	\begin{enumerate}[label=\Alph*.]
		\item Geographic distance from registered address
		\item Cash withdrawal velocity mismatch with stated charitable mission
		\item Donation from a foreign entity
		\item Large donation amount
	\end{enumerate}
	\textbf{Correct Answer: B}
	
	\subsection{3. FATF Recommendation 8 (NPOs)}
	Under FATF Recommendation 8, countries are required to apply targeted risk-based supervision of non-profit organizations. Which of the following is NOT a stated FATF focus area for NPOs vulnerable to TF?
	\begin{enumerate}[label=\Alph*.]
		\item NPOs operating without formal legal registration
		\item NPOs with large cross-border fund transfers
		\item NPOs whose expenditures do not align with their stated objectives
		\item NPOs that pay all employees via direct deposit to regulated banks
	\end{enumerate}
	\textbf{Correct Answer: D}
	
	\section{Proliferation Financing}
	
	\subsection{4. UNSCR 1540 Obligations}
	United Nations Security Council Resolution 1540 (2004) requires member states to do which of the following with respect to proliferation financing?
	\begin{enumerate}[label=\Alph*.]
		\item Establish and enforce effective controls over chemical, biological, radiological, and nuclear (CBRN) weapons-related materials
		\item Freeze all assets of any country developing nuclear weapons
		\item Prohibit all cross-border trade in dual-use goods
		\item Mandate that banks report all transactions involving any manufacturing sector
	\end{enumerate}
	\textbf{Correct Answer: A}
	
	\subsection{5. Dual-Use Goods Red Flags}
	A compliance officer at a trade finance bank reviews a transaction involving the export of precision ball bearings, industrial-grade centrifuges, and nickel alloy tubing. The destination is a front company in a third country with no legitimate metalworking or aerospace industry. The end-user is listed as a university in a country under UNSC sanctions for nuclear weapons development. What is the correct risk classification?
	\begin{enumerate}[label=\Alph*.]
		\item Low risk – commodities are standard industrial goods
		\item Proliferation financing red flag – potential nuclear weapons-related supply chain
		\item Trade-based money laundering only
		\item Sanctions evasion via real estate investment
	\end{enumerate}
	\textbf{Correct Answer: B}
	
	\subsection{6. Sectoral Sanctions vs. Proliferation Financing}
	A bank processes a wire transfer for a Russian energy company that is not individually designated on the OFAC SDN list but is subject to EU sectoral sanctions due to deep-water oil exploration activities. The funds are destined for the company’s subsidiary in a third country. Which compliance obligation applies?
	\begin{enumerate}[label=\Alph*.]
		\item No obligation because the entity is not on the SDN list
		\item Prohibition of all transactions involving any Russian entity
		\item Application of enhanced due diligence and, depending on jurisdiction, restriction on certain maturities or activities
		\item Automatic freezing of all assets under UNSCR 1373
	\end{enumerate}
	\textbf{Correct Answer: C}
	
	\section{Virtual Assets and VASPs}
	
	\subsection{7. FATF Travel Rule for Virtual Assets}
	Under the FATF Updated Guidance (2021), which of the following is required for virtual asset transfers (crypto-to-crypto or crypto-to-fiat) exceeding a threshold (typically 1,000 USD/EUR)?
	\begin{enumerate}[label=\Alph*.]
		\item No requirement – crypto is anonymous by design
		\item Originator and beneficiary information must be collected and shared with counterpart VASPs
		\item Only the transaction hash must be recorded
		\item The VASP must obtain a court order before any transfer
	\end{enumerate}
	\textbf{Correct Answer: B}
	
	\subsection{8. Anonymity-Enhanced Cryptocurrencies (AECs)}
	A customer deposits Monero (XMR) into a VASP, immediately exchanges it for Bitcoin, and then sends the Bitcoin to a non-custodial wallet. The customer refuses to provide any explanation of the source of the XMR. Under FATF standards, how should the VASP treat this activity?
	\begin{enumerate}[label=\Alph*.]
		\item As a standard transaction because Monero is legal
		\item As a high-risk indicator requiring EDD and potential SAR/STR filing
		\item As a de minimis activity below any reporting threshold
		\item As a request for simplified due diligence
	\end{enumerate}
	\textbf{Correct Answer: B}
	
	\subsection{9. Decentralized Finance (DeFi) and AML Responsibility}
	Under the FATF’s decentralized finance guidance, which entity bears AML/CFT obligations when a DeFi protocol is controlled or sufficiently decentralized but still has a “beneficial owner” or governing organization?
	\begin{enumerate}[label=\Alph*.]
		\item No entity – DeFi is unregulated globally
		\item The individual developers if they retain control or significant influence
		\item The blockchain protocol itself
		\item All users of the protocol equally
	\end{enumerate}
	\textbf{Correct Answer: B}
	
	\subsection{10. Mixers and Tumblers}
	A VASP detects that a customer sent funds through a non-regulated crypto mixing service before depositing to the VASP. The customer claims it was for “privacy protection.” What is the proper AML response?
	\begin{enumerate}[label=\Alph*.]
		\item Accept the explanation – privacy is a right
		\item Flag as high risk, apply EDD, and consider SAR filing due to obfuscation of transaction trail
		\item Only take action if the amount exceeds \$100,000
		\item Automatically close the account without investigation
	\end{enumerate}
	\textbf{Correct Answer: B}
	
	\section{Sanctions Screening and Evasion}
	
	\subsection{11. OFAC 50\% Rule}
	Company A is owned 50\% by Individual X, who is on the OFAC SDN list. Company B is owned 50\% by Company A and 50\% by Individual Y (non-sanctioned). What is the sanctions status of Company B?
	\begin{enumerate}[label=\Alph*.]
		\item Not sanctioned because Individual Y owns 50\% directly
		\item Sanctioned by application of the 50\% rule – Company A’s 50\% ownership of B makes B a blocked entity
		\item Only Company A is blocked; B remains unrestricted
		\item Requires OFAC specific license to determine
	\end{enumerate}
	\textbf{Correct Answer: B}
	
	\subsection{12. Re-routing to Evade Sanctions}
	A bank in Country A identifies that a customer is shipping goods to Country B (sanctioned) but routing payments through a shell company in Country C (non-sanctioned), with invoices showing “transshipment to Country C” but bills of lading showing final destination Country B. What typology is this?
	\begin{enumerate}[label=\Alph*.]
		\item Trade-based money laundering
		\item Sanctions evasion through transshipment
		\item Black market peso exchange
		\item Structuring of export declarations
	\end{enumerate}
	\textbf{Correct Answer: B}
	
	\subsection{13. Secondary Sanctions Risk}
	A non-U.S. bank has no U.S. presence but processes a USD clearing transaction that indirectly benefits a person on OFAC’s SDN list. The transaction does not touch the U.S. financial system. Under U.S. secondary sanctions policy, what risk does the bank face?
	\begin{enumerate}[label=\Alph*.]
		\item No risk because there is no U.S. nexus
		\item Risk of being designated on OFAC’s SDN list or losing correspondent access
		\item Only a civil penalty, not criminal
		\item The bank is automatically exempt under EU blocking statutes
	\end{enumerate}
	\textbf{Correct Answer: B}
	
	\section{KYC/CDD/EDD Nuances}
	
	\subsection{14. Beneficial Ownership of Trusts}
	A trust has the following structure: Settlor (individual A), Trustee (corporate trustee), Protector (individual B), Beneficiaries (individuals C, D, and E, each with 20\% interest, remainder discretionary). Under FATF standards, who is the ultimate beneficial owner?
	\begin{enumerate}[label=\Alph*.]
		\item Only the settlor
		\item Only the trustee
		\item Settlor, protector, and any beneficiary with >25\% interest; plus the trustee as legal owner
		\item No UBO – trusts are exempt
	\end{enumerate}
	\textbf{Correct Answer: C}
	
	\subsection{15. Legal Persons with Bearer Shares}
	A customer presents a corporate account application for a company incorporated in a jurisdiction that still allows bearer shares. The customer provides a notarized statement that the bearer shares have been immobilized with a licensed custodian. Under FATF Recommendation 24, is this sufficient for CDD?
	\begin{enumerate}[label=\Alph*.]
		\item Yes – immobilization fully mitigates risk
		\item No – jurisdictions with bearer shares are automatically unacceptable
		\item Possibly, but the bank must verify that the custodian can identify the bearer share holder at all times and that the jurisdiction prohibits future issuance
		\item Only if the customer also provides audited financial statements
	\end{enumerate}
	\textbf{Correct Answer: C}
	
	\subsection{16. Source of Funds vs. Source of Wealth}
	In Enhanced Due Diligence for a PEP, the client provides a detailed source of wealth (inheritance, business sale) but refuses to provide source of funds for a specific \$5 million wire deposit. What is the correct action?
	\begin{enumerate}
		\item Accept source of wealth as sufficient
		\item Reject the transaction because source of funds for the specific deposit is required under EDD
		\item File a SAR but process the transaction
		\item Request a waiver from the FIU
	\end{enumerate}
	\textbf{Correct Answer: B}
	
	\section{FATF Effectiveness and Mutual Evaluations}
	
	\subsection{17. Immediate Outcome 1 (IO.1)}
	In the FATF Mutual Evaluation process, Immediate Outcome 1 assesses: “Money laundering and terrorist financing risks are understood and, where appropriate, actions coordinated domestically.” Which evidence best supports a positive rating for IO.1?
	\begin{enumerate}[label=\Alph*.]
		\item A thick legal code with criminalization of ML
		\item A national risk assessment that is regularly updated and drives resource allocation
		\item High number of STRs filed
		\item Membership in the Egmont Group
	\end{enumerate}
	\textbf{Correct Answer: B}
	
	\subsection{18. Technical Compliance vs. Effectiveness}
	A country has perfect laws on paper for all 40 FATF Recommendations but has never convicted a money launderer due to lack of political will. What is the likely FATF rating?
	\begin{enumerate}[label=\Alph*.]
		\item Compliant on Technical Compliance; Low on Effectiveness (e.g., IO.7 – ML investigations)
		\item Non-compliant on Technical Compliance
		\item Partially compliant on both
		\item High on Effectiveness because laws exist
	\end{enumerate}
	\textbf{Correct Answer: A}
	
	\subsection{19. FATF Grey List Consequences}
	When a jurisdiction is placed on the FATF “Increased Monitoring” list (Grey List), what is the most immediate practical impact for financial institutions dealing with that jurisdiction?
	\begin{enumerate}[label=\Alph*.]
		\item Complete prohibition of any transactions
		\item Requirement to apply enhanced due diligence and potentially countermeasures
		\item No impact – it is only a public statement
		\item Automatic freezing of all assets from that jurisdiction
	\end{enumerate}
	\textbf{Correct Answer: B}
	
	\section{Internal Controls, Auditing, and Whistleblowers}
	
	\subsection{20. Independent Testing Scope}
	Under the FATF Standards and Wolfsberg principles, the independent audit function of an AML program must do which of the following?
	\begin{enumerate}[label=\Alph*.]
		\item Only review policies, not test transactions
		\item Test the operational effectiveness of transaction monitoring, sanctions screening, and CDD implementation
		\item Be performed by the same staff who design the AML software
		\item Report to the Chief Compliance Officer exclusively
	\end{enumerate}
	\textbf{Correct Answer: B}
	
	\subsection{21. Whistleblower Protections}
	A compliance analyst reports suspected money laundering internally. Three days later, the analyst is terminated for “performance issues.” The analyst has documented evidence of retaliation. Under which framework is this violation most directly addressed?
	\begin{enumerate}[label=\Alph*.]
		\item FATF Recommendation 33 (Protection of whistleblowers)
		\item Basel III capital requirements
		\item EU GDPR data protection rights
		\item Wolfsberg Correspondent Banking Due Diligence Questionnaire
	\end{enumerate}
	\textbf{Correct Answer: A}
	
	\subsection{22. Recordkeeping Requirements}
	Under FATF Recommendation 11, how long must financial institutions retain records on transactions and customer identification?
	\begin{enumerate}[label=\Alph*.]
		\item 1 year from transaction date
		\item At least 5 years from transaction date or account closure, whichever is later
		\item 10 years from the date of SAR filing
		\item Permanently for all accounts
	\end{enumerate}
	\textbf{Correct Answer: B}
	
	\section{Cross-Border Correspondent Banking Nuances}
	
	\subsection{23. Payable-Through Accounts (PTAs)}
	A respondent bank requests a payable-through account at a U.S. correspondent bank. The respondent bank’s customers will have direct access to the correspondent’s payment systems. Under U.S. regulations, what is the correspondent bank’s primary obligation?
	\begin{enumerate}[label=\Alph*.]
		\item No special obligation – PTA is same as any correspondent account
		\item Must apply all CDD requirements directly to the respondent bank’s customers, including identifying the underlying accountholders
		\item Only required to screen the respondent bank, not its customers
		\item Must obtain a waiver from FinCEN
	\end{enumerate}
	\textbf{Correct Answer: B}
	
	\subsection{24. Nesting Arrangements}
	A correspondent bank discovers that its respondent bank has allowed a third-party bank (not approved by the correspondent) to route payments through the correspondent account without disclosure. What is the most critical risk?
	\begin{enumerate}[label=\Alph*.]
		\item Operational inefficiency
		\item Complete loss of visibility into ultimate originators and beneficiaries (due diligence gap)
		\item Higher transaction fees
		\item Breach of data privacy laws
	\end{enumerate}
	\textbf{Correct Answer: B}
	
	\section{Money Mules and Human Trafficking Financial Indicators}
	
	\subsection{25. Money Mule Recruitment Typology}
	An account shows multiple incoming peer-to-peer transfers from dozens of different individuals under \$500 each, followed by immediate aggregated wire transfers to an offshore shell company account. The account holder is a college student with no verifiable income. What is the most likely typology?
	\begin{enumerate}[label=\Alph*.]
		\item Personal savings consolidation
		\item Money mule account for layering stolen funds
		\item Charitable donation collection
		\item Cryptocurrency mining proceeds
	\end{enumerate}
	\textbf{Correct Answer: B}
	
	\subsection{26. Human Trafficking Financial Red Flags}
	Which of the following transaction patterns is most indicative of human trafficking for forced labor rather than smuggling?
	\begin{enumerate}[label=\Alph*.]
		\item Single large wire transfer from a family to a coyote
		\item Multiple small, regular payments from a corporate account to several individuals with the same surname, no employment contract on file
		\item One-time cash deposit at an ATM
		\item Purchase of airline tickets for a family vacation
	\end{enumerate}
	\textbf{Correct Answer: B}
	
	\section{Conclusion – Mixed High-Difficulty}
	
	\subsection{27. Reverse Correspondent Banking Risk}
	A small bank in a developing country maintains a correspondent account for a larger international bank (reverse correspondent). The international bank processes USD clearing through the small bank’s account. What unique risk arises?
	\begin{enumerate}[label=\Alph*.]
		\item No unique risk – reverse correspondent is safer
		\item The small bank may be used as a blind gateway into the larger bank’s client base
		\item The larger bank must close the account
		\item Only foreign exchange risk exists
	\end{enumerate}
	\textbf{Correct Answer: B}
	
	\subsection{28. Shell Bank Prohibition – USA PATRIOT Act Section 313}
	Under Section 313 of the USA PATRIOT Act, a “shell bank” is defined as:
	\begin{enumerate}[label=\Alph*.]
		\item Any bank with less than \$100 million in assets
		\item A bank that has no physical presence in any country and is not affiliated with a regulated financial institution
		\item Any bank operating from a tax haven
		\item A bank that only serves offshore companies
	\end{enumerate}
	\textbf{Correct Answer: B}
	
	\subsection{29. Currency Transaction Reports (CTRs) and Structuring}
	A customer deposits \$9,000 cash on Monday, \$9,000 on Wednesday, and \$9,000 on Friday. The bank files a CTR for the aggregate \$27,000. However, the customer later sues the bank for violating his privacy. Under U.S. law, is the bank protected?
	\begin{enumerate}[label=\Alph*.]
		\item No – CTR filing requires customer consent
		\item Yes – banks are required to file CTRs for aggregate cash transactions > \$10,000 in a single business day, and structuring itself is illegal
		\item Only if the customer is a PEP
		\item No – the bank should have filed a SAR instead
	\end{enumerate}
	\textbf{Correct Answer: B}
	
	\subsection{30. FATF Recommendation 16 (Wire Transfer Rule)}
	Under FATF Recommendation 16, which of the following is NOT required information to accompany a cross-border wire transfer?
	\begin{enumerate}[label=\Alph*.]
		\item Originator’s name
		\item Originator’s account number
		\item Originator’s country of birth
		\item Originator’s address or national ID number or date of birth
	\end{enumerate}
	\textbf{Correct Answer: C}
	
\end{document}